\section{Escapes at both ends of a random path}\label{sec:escapes}

In this section we prove \cref{prop:escapes}. We assume \cref{ass:G,ass:U}, fix $p\in[p_c/2,p_c)$ and
write $T=T_p$. The construction is the one of~\cite{Ope26}, with one simplification available in the
transitive case: the geodesic segments below are chosen among all geodesic segments, rather than within
the orbit of a fixed one.

We must find, for a positive fraction of lazy-walk paths $w$, short paths from both endpoints that avoid
the interior of the corridor and end at vertices from which $T$ sees little of the neighbourhood of the
corridor. Two observations make this possible. First, after a few steps an independent random walk
averages the connection kernel, and this averaging does not depend on $\norm T$
(\cref{lem:smoothing}). Second, if the raw path starts and ends with halves of geodesic segments, an
escape can leave along the other half of the segment and is then unlikely to meet the raw path, by the
decay of $P^l$ on $\ell^2(V)$. A symmetry of the construction under reversal, obtained from the
mass-transport principle, transfers the estimate from one endpoint to the other.

\begin{lemma}[Smoothing]\label{lem:smoothing}
For all $a,b\in V$ and $l\ge0$, $0\le(P^lT)(a,b)\le\rho^l$.
\end{lemma}

\begin{proof}
Since $T$ is positive semidefinite (\cref{lem:subcritical}), the Cauchy--Schwarz inequality for the form
$(f,g)\mapsto\langle f,Tg\rangle$ gives
$(P^lT)(a,b)^2=\langle P^l\1_a,T\1_b\rangle^2\le(P^lTP^l)(a,a)\,T(b,b)$, and $T(b,b)=1$. The
operator $P^lTP^l$ is invariant and positive semidefinite, so by \cref{lem:trace},
\begin{equation*}
  (P^lTP^l)(a,a)=\tr(P^lTP^l)=\tr(P^{2l}T)\le\norm{P^{2l}}\tr(T)=\rho^{2l}\,,
\end{equation*}
since $\tr(T)=T(o,o)=1$.
\end{proof}

\paragraph{Geodesic segments.}
For $J\in\N$ let $\mathsf S_J$ be the set of geodesic paths $g=(g_{-J},\dots,g_J)$ in~$G$, that is, of
sequences of vertices with $\dist(g_i,g_{i'})=\abs{i-i'}$. Since $G$ is infinite, connected and locally
finite, there are vertices at distance $2J$, so $\mathsf S_J\neq\emptyset$; it is invariant under
automorphisms. For $x\in V$ and $\abs i\le J$, the number $n_i\coloneqq\#\{g\in\mathsf S_J\colon
g_i=x\}$ is finite, because such $g$ lie in $B(x,2J)$, and positive and independent of $x$ by
transitivity. The mass-transport principle applied to $\#\{g\in\mathsf S_J\colon g_i=x,\,g_{i'}=y\}$
shows that $n_i=n_{i'}$ for all $i,i'$; we write $n_\star$ for this number. Consequently, for
$\abs i,\abs{i'}\le J$ the \emph{position kernel}
\begin{equation*}
  \Pi_{i\to i'}(x,y)\coloneqq\frac{\#\{g\in\mathsf S_J\colon g_i=x,\ g_{i'}=y\}}{n_\star}
\end{equation*}
is doubly stochastic, so $\norm{\Pi_{i\to i'}}\le1$.

\paragraph{The labelled experiment.}
Let $k>2J$ and $m\in\N$. A \emph{labelled object} rooted at $u$ is a tuple $\Xi=(g,\pi,g',\zeta,\zeta')$
consisting of $g\in\mathsf S_J$ with $g_0=u$; a lazy-walk path $\pi=(\pi_0,\dots,\pi_{k-2J})$ with
$\pi_0=g_J$; a segment $g'\in\mathsf S_J$ with $g'_J=\pi_{k-2J}$; and lazy-walk paths $\zeta$ and
$\zeta'$ of length $m$ starting at $g_{-J}$ and $g'_{-J}$. Under the probability measure
\begin{equation}\label{eq:labelled-law}
  \mathbb Q_u(\Xi)\coloneqq\frac{\mu(\pi)\,\mu(\zeta)\,\mu(\zeta')}{n_\star^2}
\end{equation}
the segment $g$ is uniform, $\pi$ is a lazy walk from $g_J$, the segment $g'$ is uniform given its
position $g'_J$, and $\zeta,\zeta'$ are independent lazy walks. The object determines (see
\cref{fig:escapes}):
\begin{itemize}
  \item the \emph{raw path}, a $k$-step lazy-walk path with base $u$ and end $v\coloneqq g'_0$,
  \begin{equation*}
    w(\Xi)\coloneqq(g_0,g_1,\dots,g_J,\pi_1,\dots,\pi_{k-2J},g'_{J-1},\dots,g'_0)\,;
  \end{equation*}
  \item the \emph{escapes} from $u$ and from~$v$, with endpoints $Z\coloneqq\zeta_m$ and
  $Z'\coloneqq\zeta'_m$,
  \begin{equation*}
    \psi\coloneqq(g_0,g_{-1},\dots,g_{-J},\zeta_1,\dots,\zeta_m)\,,\qquad
    \psi'\coloneqq(g'_0,g'_{-1},\dots,g'_{-J},\zeta'_1,\dots,\zeta'_m)\,.
  \end{equation*}
\end{itemize}

\begin{figure}[t]
\centering
% The labelled experiment of Section 10 (schematic)
\begin{tikzpicture}[scale=0.88,
    vtx/.style={circle, fill, inner sep=1.3pt},
    raw/.style={figmain, line width=1.4pt},
    rawwalk/.style={figmain, line width=1.4pt, decorate, decoration={snake, amplitude=1.1pt, segment length=9pt}},
    esc/.style={figaux, thick, dashed},
    escwalk/.style={figaux, thick, dashed, decorate, decoration={snake, amplitude=1.1pt, segment length=7pt}}]
  \coordinate (u) at (0,0);
  \coordinate (gJ) at (1.8,1.6);
  \coordinate (gmJ) at (-1.8,-1.6);
  \coordinate (hJ) at (8.2,1.6);
  \coordinate (v) at (10,0);
  \coordinate (hmJ) at (11.8,-1.6);
  \coordinate (z) at (-0.6,-3.1);
  \coordinate (zp) at (10.6,-3.1);
  % raw path: first half of g, middle walk, reversed first half of g'
  \draw[raw] (u) -- (gJ);
  \draw[rawwalk] (gJ) .. controls (4,2.4) and (6,0.8) .. (hJ);
  \draw[raw] (hJ) -- (v);
  % escapes: other halves of the segments, then walk tails
  \draw[esc] (u) -- (gmJ);
  \draw[escwalk] (gmJ) .. controls (-2.2,-2.6) and (-1.4,-3.4) .. (z);
  \draw[esc] (v) -- (hmJ);
  \draw[escwalk] (hmJ) .. controls (12.2,-2.6) and (11.4,-3.4) .. (zp);
  \foreach \p in {u,gJ,gmJ,hJ,v,hmJ,z,zp} \node[vtx] at (\p) {};
  \node[font=\small, anchor=east] at (u) {$u=g_0\ $};
  \node[font=\small, anchor=south east] at (gJ) {$g_J$};
  \node[font=\small, anchor=east] at (gmJ) {$g_{-J}$};
  \node[font=\small, anchor=south west] at (hJ) {$g'_J$};
  \node[font=\small, anchor=west] at (v) {$\ v=g'_0$};
  \node[font=\small, anchor=west] at (hmJ) {$g'_{-J}$};
  \node[font=\small, anchor=north] at (z) {$Z$};
  \node[font=\small, anchor=north] at (zp) {$Z'$};
  \node[font=\small, text=figmain] at (5,2.45) {middle walk $\pi$, $k-2J$ steps};
  \node[font=\small, text=figaux] at (-2.9,-2.7) {tail $\zeta$};
  \node[font=\small, text=figaux] at (12.9,-2.7) {tail $\zeta'$};
\end{tikzpicture}

\caption{The labelled experiment (schematic). The raw path (solid) consists of the positive half of a
geodesic segment $g$ through~$u$, a middle walk~$\pi$, and the positive half of a geodesic segment $g'$
through~$v$, traversed backwards. The escapes (dashed) follow the negative halves of the same segments
and continue as independent walks $\zeta,\zeta'$ of length~$m$. Reversing the middle walk and exchanging
the roles of $(g,\zeta)$ and $(g',\zeta')$ gives the experiment rooted at~$v$.}\label{fig:escapes}
\end{figure}

\paragraph{Reversal.}
The \emph{reversal} of $\Xi$ is $\Xi^\dagger\coloneqq(g',\overleftarrow\pi,g,\zeta',\zeta)$, where
$\overleftarrow\pi$ is $\pi$ traversed backwards; it is a labelled object rooted at~$v$, its raw path is
$w(\Xi)$ traversed backwards, and its two escapes are those of $\Xi$ in exchanged order. By symmetry of
$P$, $\mu(\overleftarrow\pi)=\mu(\pi)$, so $\mathbb Q_v(\Xi^\dagger)=\mathbb Q_u(\Xi)$. We claim that
\begin{equation}\label{eq:reversal}
  \EE_{\mathbb Q_o}[F]=\EE_{\mathbb Q_o}[F\circ\dagger]
\end{equation}
for every nonnegative function $F$ of labelled objects that is invariant under automorphisms. Indeed,
the mass $M_F(x,y)\coloneqq\sum\mathbb Q_x(\Xi)F(\Xi)$, summed over the objects rooted at $x$ whose raw
path ends at~$y$, is diagonally invariant. The mass sent out by $o$ is $\EE_{\mathbb Q_o}F$. Since
$\Xi\mapsto\Xi^\dagger$ is a bijection between the objects rooted at $x$ ending at $o$ and the objects
rooted at $o$ ending at~$x$, the mass received by $o$ is $\EE_{\mathbb Q_o}[F\circ\dagger]$. The
mass-transport principle gives~\eqref{eq:reversal}.

\paragraph{Comparison with the lazy walk.}
For a raw path $w$ with base $o$, we claim
\begin{equation}\label{eq:domination}
  \mathbb Q_o\bigl(w(\Xi)=w\bigr)\le(2d)^{2J}\mu(w)\,.
\end{equation}
Indeed, $w(\Xi)=w$ fixes $\pi$, the positive half of $g$ and the reversed positive half of~$g'$. Summing
\eqref{eq:labelled-law} over the at most $n_\star$ compatible segments $g$, the at most $n_\star$ compatible
segments $g'$, and all tails gives at most $\mu(\pi)$. The $2J$ steps of $w$ along the two
half-segments are steps between distinct neighbours, each of probability $1/(2d)$, so
$\mu(w)=(2d)^{-2J}\mu(\pi)$.

\begin{proof}[Proof of \cref{prop:escapes}]
Let $\lambda\coloneqq\log(1/\rho)>0$. Throughout, $\Xi$ has law $\mathbb Q_o$, so $u=o$.

\emph{Stage 1: the escape from $u$ avoids the raw path.} We bound the probability that $\psi$ visits a
vertex of $w(\Xi)$ other than~$u$; revisits of $u$ are allowed. Condition on~$g$. The deterministic part
$(g_0,\dots,g_{-J})$ of $\psi$ meets the first block $(g_0,\dots,g_J)$ of the raw path only in $u$,
because $g$ is a geodesic. We use $P^r(a,b)=\langle\1_a,P^r\1_b\rangle\le\rho^r$ throughout.
\begin{itemize}
  \item The middle walk visits a given vertex $g_{-i}$, $0\le i\le J$, at time $r$ with probability at
  most $\rho^r$, and with probability zero for $r<J$.
  \item The tail $\zeta$ visits a given vertex $g_i$, $0\le i\le J$, at time $l$ with probability at most
  $\rho^l$, and with probability zero for $l<J$.
  \item Since $\zeta$ and $\pi$ are independent given $g$, the probability that $\zeta_l=\pi_r$ is
  \begin{equation*}
    \sum_xP^l(g_{-J},x)\,P^r(g_J,x)=P^{l+r}(g_{-J},g_J)\le\rho^{l+r}\,,
  \end{equation*}
  and it is zero for $l+r<2J$.
\end{itemize}
These collisions have total probability at most
\begin{equation*}
  E_J\coloneqq2(J+1)\sum_{r\ge J}\rho^r+\sum_{s\ge2J}(s+1)\rho^s\,,
\end{equation*}
which tends to zero as $J\to\infty$; the time sums run to infinity, so the bound is uniform in $k$
and~$m$. It remains to consider collisions with the last block $(g'_J,\dots,g'_0)$. Given $g$, the
endpoint of the middle walk has law $\varphi\coloneqq P^{k-2J}(g_J,\cdot)$ with
$\norm\varphi_2\le\rho^{k-2J}$, and $g'_i$ has law $\varphi\Pi_{J\to i}$, whose $\ell^2$ norm is also at
most $\rho^{k-2J}$. The escape $\psi$ is independent of $(\pi,g')$ given $g$. A deterministic position
of $\psi$ is hit by $g'_i$ with probability at most $\norm{\varphi\Pi_{J\to i}}_2\le\rho^{k-2J}$, and the
tail position $\zeta_l$, whose law has $\ell^2$ norm at most $\rho^l$, with probability at most
$\rho^{l+k-2J}$ by the Cauchy--Schwarz inequality. A union bound over $0\le i\le J$, the $J+1$
deterministic positions and $l\ge0$ bounds these collisions by $F_J\rho^{k-2J}$ with
$F_J\coloneqq(J+1)(J+1+(1-\rho)^{-1})$. In the same way,
$\mathbb Q_o(\dist(u,v)<3)\le\abs{B(o,2)}^{1/2}\rho^{k-2J}$.

\emph{Stage 2: the escape from $v$.} The event that $\psi'$ visits a vertex of $w(\Xi)$ other than~$v$ is
the event that the escape from the root of $\Xi^\dagger$ visits a vertex of the raw path of $\Xi^\dagger$
other than the root. By~\eqref{eq:reversal} it has the same probability as the event of Stage~1.

\emph{Stage 3: choice of $J$ and $k_0$.} Fix $J$ with $2E_J\le1/8$, and then $L_0$ with
$(2F_J+\abs{B(o,2)}^{1/2})\rho^{L_0}\le1/8$. For $k-2J\ge L_0$ and every~$m$, with
$\mathbb Q_o$-probability at least $3/4$, both escapes avoid every vertex of the raw path other than
their starting points, and $\dist(u,v)\ge3$. The escapes need not avoid each other.

\emph{Stage 4: the endpoints of the escapes.} Fix $C_1$ with $\lambda C_1>1$, and let
$m\coloneqq\lceil C_1\log((k+1)R)\rceil$. Given $g,\pi,g'$, and hence the raw path, the tails are
independent lazy walks, and $\abs{\Nb[I_w]}\le(d+1)(k+1)$. By \cref{lem:smoothing},
\begin{equation*}
  \EE_{\mathbb Q_o}\Bigl[\sum_{b\in\Nb[I_{w(\Xi)}]}T(Z,b)\Bigm|g,\pi,g'\Bigr]
  =\sum_{b\in\Nb[I_{w(\Xi)}]}(P^mT)(g_{-J},b)\le(d+1)(k+1)\rho^m\,,
\end{equation*}
and the same holds for~$Z'$. By Markov's inequality and $\rho^m\le((k+1)R)^{-\lambda C_1}$, the
probability that one of the two sums exceeds $1/(4R)$ is at most
$8(d+1)((k+1)R)^{1-\lambda C_1}$. Fix $k_0\ge2J+L_0$ such that this is at most $1/4$ for all
$k\ge k_0$ and $R\ge1$.

\emph{Stage 5: conclusion.} For $k\ge k_0$, with $\mathbb Q_o$-probability at least $1/2$ the events of
Stages~3 and~4 occur together. Then $\abs{\gamma_w}\ge\dist(u,v)\ge3$ for $w=w(\Xi)$, and the interior
$I_w$ consists of vertices of the raw path other than $u$ and~$v$. So $\psi$ avoids $I_w$, and erasing its
holding steps and loops yields a self-avoiding path in $\Lambda_w$ from $u$ to $Z$ of length at most
$J+m\le j$; similarly for~$v$. Hence $w(\Xi)\in\mathcal W$ on this event, and~\eqref{eq:domination}
gives
\begin{equation*}
  \Erw\bigl[\1\{w\in\mathcal W\}\bigr]\ge(2d)^{-2J}\,\mathbb Q_o\bigl(w(\Xi)\in\mathcal W\bigr)
  \ge\tfrac12(2d)^{-2J}\eqqcolon c_\star\,.
\end{equation*}
The constants $J$, $k_0$, $C_1$ and $c_\star$ depend only on~$G$.
\end{proof}
